% Generated draft: requires mathematical and source review.
\documentclass[11pt]{article}
\usepackage[margin=25mm]{geometry}
\usepackage{fontspec}
\setmainfont{Latin Modern Roman}
\usepackage{amsmath,amssymb,mathtools,unicode-math}
\setmathfont{Latin Modern Math}
\usepackage{enumitem,needspace,xurl,hyperref}
\setlength{\parindent}{0pt}
\setlength{\parskip}{5pt}
\setlength{\emergencystretch}{3em}
\newcommand{\N}{\mathbb{N}}
\newcommand{\Z}{\mathbb{Z}}
\newcommand{\Q}{\mathbb{Q}}
\newcommand{\R}{\mathbb{R}}
\newcommand{\C}{\mathbb{C}}
\newcommand{\F}{\mathbb{F}}
\newcommand{\A}{\mathbb{A}}
\newcommand{\PP}{\mathbb{P}}
\newcommand{\E}{\mathbb{E}}
\newcommand{\eps}{\varepsilon}
\newcommand{\dd}{\,\mathrm{d}}
\begin{document}
\section*{UnsolvedMath ID: 700349}
\textbf{Source problem code:} catalogue entry
\subsection{Definitions and mathematical statement}
Let \(G\) be an arbitrary group, with no finiteness assumption. The supplied record does not define the notation \(C\) or the term \(C\)-closed. The following formulation uses the usual centralizer interpretation; this convention requires confirmation against the cited source.

For a subset \(S\subseteq G\), define its centralizer by
\[
C_G(S)=\{g\in G:gs=sg\text{ for every }s\in S\}.
\]
Under this interpretation, a subgroup \(H\leq G\) is \(C\)-closed if
\[
H=C_G(C_G(H)).
\]
Equivalently, \(H\) is the centralizer of some subset of \(G\). Write
\[
\mathcal C(G)=\{H\leq G:H=C_G(C_G(H))\}.
\]
Ordered by inclusion, this is a complete lattice. Its binary operations, as given in the supplied statement, are
\[
A\wedge_{\mathcal C(G)}B=A\cap B,
\qquad
A\vee_{\mathcal C(G)}B=C_G\bigl(C_G(A)\cap C_G(B)\bigr).
\]
Let \(\operatorname{Sub}(G)\) denote the lattice of all subgroups of \(G\), whose binary meet and join are intersection and generated subgroup, respectively.

\textbf{Problem.} Describe all groups \(G\) for which \(\mathcal C(G)\) is a sublattice of \(\operatorname{Sub}(G)\). Explicitly, classify the groups satisfying
\[
\forall A,B\in\mathcal C(G),\qquad
\langle A\cup B\rangle
=C_G\bigl(C_G(A)\cap C_G(B)\bigr).
\]
Equivalently, the subgroup generated by any two \(C\)-closed subgroups must again be \(C\)-closed. Here sublattice means preservation of binary meets and joins; preservation of arbitrary or empty joins in the ambient subgroup lattice is not imposed.
\subsection{Short English statement}
Using the usual centralizer meaning of \(C\), classify all groups in which the subgroup generated by any two double-centralizer-closed subgroups is again double-centralizer-closed. The source does not explicitly define this terminology, so that interpretation requires review. The registry labels the problem open; this is supplied status information, not an independently verified assessment.
\subsection{Sources}
\begin{itemize}
\item Supplied local UnsolvedMath registry record, permanent ID \(700349\), exact source code \texttt{catalogue entry}, titled \emph{Kourovka Notebook Problem 6.2}. The record supplies the status \texttt{open}.
\item \emph{The Kourovka Notebook (Unsolved Problems in Group Theory)}, problem 6.2, originating in the 6th issue (1978); the supplied background identifies the 21st issue edition (2026), \texttt{arXiv:1401.0300v48}. This bibliographic attribution is supplied evidence and has not been independently checked.
\item Proposed by V. A. Antonov and N. F. Sesekin, with proposed year 1978. The supplied collection description identifies E. I. Khukhro and V. D. Mazurov as editors of the Issues 1--21 collection (1965--2026).
\end{itemize}
\end{document}
